\documentclass[a4paper,twoside]{article}


% --- RUIMTE EN LAYOUT ---
\usepackage{setspace}

\usepackage{geometry} %size regulating
\geometry{
 a4paper,
 left=25mm,
 right=25mm,
 top=20mm,
 bottom=20mm,
 }

\usepackage{parskip} % for better spacing between paragraphs
\usepackage{flushend} % Alleen nodig als je laatste pagina van artikel wilt vullen met tekst, zodat de kolommen even lang zijn. Anders kan je dit beter weglaten.

% --- WISKUNDE ENSYMBOLEN ---
\usepackage{amsmath, amssymb, amsfonts, amsthm} % For nicer math format
\usepackage{mathrsfs} % For curly arrr
\usepackage{bm} % For bold symbols eg \bm{\nabla}
\usepackage{esint} % For bolletje in integrals
\usepackage{dsfont}
\usepackage{siunitx} % Voeg deze expliciet toe voor die graden en eenheden!
\usepackage{braket} % For bra-ket notation
\usepackage{mathtools}
\usepackage{url}
\allowdisplaybreaks

% --- OVERIGE ---

\usepackage{graphicx} % Required for inserting images
\usepackage[hidelinks]{hyperref} % geen kaders om linken
\usepackage{subcaption}
%\usepackage[dutch]{babel}
\usepackage[T1]{fontenc}
\usepackage[utf8]{inputenc}
\usepackage{lmodern}
\usepackage{multicol}


% --- Tikz packages ---
\usepackage{tikz}
\usetikzlibrary{angles,quotes, babel} % for pic (angle labels)
\usetikzlibrary{arrows.meta,decorations.markings,calc, decorations.pathmorphing, decorations.pathreplacing, positioning, shapes, shadings, patterns}
\usepackage{xcolor}
\usepackage{tikz-3dplot}
\usepackage[siunitx]{circuitikz}

\usepackage{etoolbox} % ifthen
\usepackage[outline]{contour} % glow around text
\usepackage{xfp} % higher precision (16 digits?)
\contourlength{1.1pt}



% --- THEOREM STIJLEN MET EXTRA RUIMTE ---

\theoremstyle{definition}
\newtheorem{postulate}{Postulaat}
\newtheorem{definition}{Definitie}
\newtheorem{derivation}{Afleiding}
\newtheorem{proposition}{Propositie}
\newtheorem{example}{Voorbeeld}

\theoremstyle{plain}
\newtheorem{theorem}{Theorem}


% --- EXTRA COMMANDO'S ---

\newcommand{\dbar}{{\mkern+3mu\mathchar'26\mkern-12mu d}}

% --- Titel en Introductie --- 

\title{Kwantum II: Moeilijkheden}
\author{T. Cools}
\date{2026-2027}
\raggedbottom

% Toelaten dat pagina-overgangen plaatsvinden tussen opeenvolgende (lege) titels
\makeatletter
\let\old@startsection\@startsection
\renewcommand{\@startsection}[6]{%
  \@nobreakfalse
  \old@startsection{#1}{#2}{#3}{#4}{#5}{#6}%
}
\makeatother

\begin{document}
\maketitle

\begin{abstract}
   Hier noteer ik de delen van de cursus die ik nog niet helemaal begrijp of moeilijk vind.
\end{abstract}

\section{Formele ontwikkeling van de kwantummechanica}
\subsection{Inleiding}
\subsection{De elektronspin en het superpositiebeginsel}
\subsubsection{Spinkwantisatie}
\subsubsection{Experiment 1: SGz na SGz}
\subsubsection{Experiment 2: SGx na SGz}
\subsubsection{Experiment 3: SGz na SGx na SGz}
\subsubsection{Het superpositiebeginsel}
\subsection{Elektronspin en kwantummechanica in 2D Hilbertruimte}
\subsubsection{Het superpositiebeginsel, vectorruimten en waarschijnlijkheden}
\subsubsection{Experiment 4: SGx na SGx na SGz en de collaps van de toestandsvector}
\subsubsection{Superpositiebeginsel, Hilbertruimte en waarschijnlijkheden}
\subsubsection{Hermitische operatoren, matrixelementen en verwachtingswaarden van observabelen}
\subsubsection{Het superpositiebeginsel en interferentie}
\subsection{De postulaten van de kwantummechanica}
\subsection{Hilbertruimte, klassieke faseruimte en het correspondentieprincipe}
\subsection{Reële lineaire algebra in een notendop}
\subsection{Complexe lineaire algebra}
\subsubsection{Hilbertruimte en bra- en ket-notatie, resolutie van de identiteit}
\subsubsection{Lineaire operatoren}
\subsubsection{Hermitische operatoren}
Theorema 6 bewijs met jordaanblokken!
\subsection{Veralgemening tot een oneindig-dimensionale Hilbertruimte}
\subsubsection{Plaatseigenkets en plaatsmetingen}
\subsubsection{Intermezzo over de Dirac-delta distributie}
\subsubsection{De Hilbertruimte van de kwadratisch integreerbare functies}
\subsubsection{Ontwikkeling van een toestand in een basis van eigenfuncties van een algemene operator $\hat{A}$}
\subsection{Het direct product van Hilbertruimten}
\subsection{Unitaire transformaties}
\subsubsection{Definities en eigenschappen}
\subsubsection{Verandering van basis en representaties}
\subsubsection{Infinitesimale unitaire transformaties}
\subsubsection{De evolutie-operator}
\subsubsection{Het Heisenbergbeeld}

\section{Symmetrieën en behoudswetten in kwantummechanica}
\subsection{Symmetrieën, behoudswetten en degeneratie}
\subsubsection{Symmetrieën en behoudswetten in klassieke mechanica}
\subsubsection{Symmetrieën en behoudswetten in kwantummechanica}
\subsubsection{Symmetrieën en degeneratie van het energiespectrum}
\subsection{Translatiesymmetrie}
\subsubsection{Infinitesimale translaties en de momentumoperator}
\subsubsection{De momentumoperator in de configuratierepresentatie}
\subsubsection{De impulsrepresentatie}
\subsubsection{De onbepaaldheidsrelatie van Heisenberg}
\subsection{Rotatiesymmetrie}
\subsubsection{Infinitesimale rotaties in kwantummechanica}
\subsubsection{Spinimpulsmoment en totaal draaimoment}
\subsubsection{Eigenvectoren en eigenwaarden van het draaimoment}
\subsubsection{Matrixelementen van draaimomentoperatoren en representaties van de rotatiegroep}
\subsubsection{Voorbeelden van representaties van de rotatiegroep}
\subsubsection{Samenstelling van draaimoment}
\subsubsection{O(3) en SU(2)}
\subsubsection{Het baanimpulsmoment in de configuratierepresentatie}
\subsubsection{Eigenschappen van sferische harmonieken}
\subsubsection{Sferische harmonieken met behulp van ladderoperatoren}

\section{Driedimensionale vraagstukken in de kwantummechanica}
\subsection{De driedimensionale Schrödingervergelijking}
\subsection{Het vrij deeltje in Cartesische coördinaten: vlakke golven}
\subsection{Een deeltje in een box}
\subsection{Een deeltje in een centrale potentiaal}
\subsubsection{Het vrij deeltje in sferische coördinaten}
\subsection{De ontwikkeling van een vlakke golf in sferische harmonieken}
\subsection{Het waterstofatoom}
\subsubsection{Schrödingervergelijking voor een tweedeeltjessysteem}
\subsubsection{Energieniveaus van het waterstofatoom}
\subsubsection{Golffuncties van de gebonden toestanden}
\subsection{Schrödingervergelijking van een deeltje in een elektromagnetisch veld}
\subsection{Het ijkprincipe}

\section{Verstrooiingstheorie}
\subsection{Verstrooiingsexperimenten en werkzame doorsneden}
\subsection{Potentiaalverstrooiing, randvoorwaarden en Greense functies}
\subsection{Lippmann-Schwinger verstrooiingstheorie}
\subsection{De Bornapproximatie voor de Yukawapotentiaal}
\subsection{Partiële golfanalyse}
\subsection{Toepassing van de partiële golfmethode}
\subsubsection{Verstrooiing door een sferische vierkante potentiaalput}
\subsubsection{Verstrooiing aan een harde sfeer}

\section{Benaderingstechnieken}
\subsection{Perturbatietheorie voor een niet-ontaard energieniveau}
\subsection{Tijdsonafhankelijke perturbatietheorie voor een ontaard energieniveau}
\subsection{Toepassing: de fijnstructuur van het waterstofatoom}
\subsubsection{Energieverschuiving t.g.v.\ $\hat{V}_1'$ (relativistische correctie tot de kinetische energie)}
\subsubsection{Energieverschuiving t.g.v.\ $\hat{V}_2'$ (spin-baankoppeling)}
\subsubsection{Energieverschuiving t.g.v.\ $\hat{V}_3'$ (Darwinterm)}
\subsubsection{Totale energieverschuiving: fijnstructuur van het waterstofatoom}
\subsection{Tijdsafhankelijke perturbatietheorie}
\subsubsection{Tijdsafhankelijke potentialen en het interactiebeeld}
\subsubsection{Tijdsafhankelijke perturbatietheorie: de Dysonreeks}
\subsubsection{De transitiewaarschijnlijkheid}
\subsubsection{Een constante perturbatie}
\subsubsection{Een harmonische perturbatie}
\subsection{Toepassing: interactie met een klassiek stralingsveld}
\subsection{De variationele methode}
\subsection{Toepassing van de variationele methode: de anharmonische oscillator}

\section{Relativistische kwantummechanica: de Diracvergelijking}
\subsection{Vroegste pogingen tot een relativistische kwantummechanica}
\subsection{De Diracvergelijking}
\subsection{Niet-relativistische limiet}
\subsection{Covariante vorm van de Diracvergelijking}
\subsection{Lorentzcovariantie van de Diracvergelijking}
\subsection{Pariteitstransformatie}
\subsection{Bilineaire covariante vormen}
\subsection{Oplossing van de Diracvergelijking voor een vrij deeltje}
\subsection{Golfpakketten}
\subsection{De gatentheorie van Dirac}

\section{Tweede kwantisatie}
\subsection{Toestanden met twee deeltjes}
\subsection{Toestanden met $N$ deeltjes}
\subsubsection{Fermionen}
\subsubsection{Bosonen}
\subsection{Creatieoperatoren en annihilatieoperatoren}
\subsubsection{Fermion creatie- en annihilatieoperatoren}
\subsubsection{Boson creatie- en annihilatieoperatoren}
\subsection{Operatoren in tweede kwantisatie}
\subsubsection{Eéndeeltje operatoren}
\subsubsection{Diagonale vorm van een hermitische ééndeeltje operator}
\subsubsection{Tweedeeltjesoperatoren}




\bibliographystyle{plain}
\bibliography{/Users/tristancools/Documents/Overleaf-Latex-Typst/References/references}

\end{document}